% draft-sigurdsson-t32-00: T32 as an Internet-Draft.
%
% Laid out like an IETF Internet-Draft (header block, boilerplate, running
% header and footer, BCP 14 language, Security and IANA sections, split
% references), following the structure of RFC 9562. Its normative content is
% draft 0.3 of SPEC.md; where the two differ, one of them has a defect.
%
% The IETF accepts submissions only as RFCXML; see README.md in this directory.

\documentclass[11pt]{article}

\usepackage[T1]{fontenc}
\usepackage[utf8]{inputenc}
\usepackage{lmodern}
\usepackage{textcomp}
\usepackage{amsmath}
\usepackage[a4paper,margin=25mm,headheight=14pt,headsep=8mm,footskip=12mm]{geometry}
\usepackage{microtype}
\usepackage{fancyhdr}
\usepackage{array,booktabs,longtable,tabularx}
\usepackage{fancyvrb}
\usepackage{enumitem}
\usepackage{xcolor}
\usepackage{url}
\def\UrlBreaks{\do\/\do.\do-}
\usepackage[hidelinks]{hyperref}

% ------------------------------------------------------------------ draft metadata
\newcommand{\draftname}{draft-sigurdsson-t32-00}
\newcommand{\shorttitle}{T32}
\newcommand{\published}{1 November 2026}
\newcommand{\pubmonth}{November 2026}
\newcommand{\expires}{5 May 2027} % 185 days after posting
\newcommand{\authorshort}{Sigurdsson}

% ------------------------------------------------------------------ house style
\setlength{\parindent}{0pt}
\setlength{\parskip}{0.6\baselineskip}
% Lets a line stretch a little before it overruns; engines differ slightly in font metrics.
\setlength{\emergencystretch}{1.5em}
\setlist{nosep,leftmargin=1.5em,topsep=0.2\baselineskip}
\renewcommand{\arraystretch}{1.15}
\setcounter{tocdepth}{2}
\setcounter{secnumdepth}{3}

% RFCs number headings "1.", "1.1.", and so on, and list them with dotted leaders.
\makeatletter
\renewcommand{\@seccntformat}[1]{\csname the#1\endcsname.\quad}
\renewcommand{\numberline}[1]{\hb@xt@\@tempdima{#1.\hfil}}
\renewcommand*{\l@section}{\@dottedtocline{1}{0em}{2.4em}}
\renewcommand*{\l@subsection}{\@dottedtocline{2}{2.4em}{2.8em}}
\makeatother

\pagestyle{fancy}
\fancyhf{}
\fancyhead[L]{Internet-Draft}
\fancyhead[C]{\shorttitle}
\fancyhead[R]{\pubmonth}
\fancyfoot[L]{\authorshort}
\fancyfoot[C]{Expires \expires}
\fancyfoot[R]{[Page \thepage]}
\renewcommand{\headrulewidth}{0pt}
\fancypagestyle{plain}{\fancyhf{}\fancyfoot[L]{\authorshort}\fancyfoot[C]{Expires \expires}\fancyfoot[R]{[Page \thepage]}}

% Stamps, durations and other literal text.
\newcommand{\lit}[1]{\texttt{#1}}
\newcommand{\dAT}{\ensuremath{\Delta\mathrm{AT}}}
\newcommand{\pow}[2]{\ensuremath{#1^{#2}}}
\newcommand{\Sec}[1]{Section~\ref{#1}}
\newcommand{\App}[1]{Appendix~\ref{#1}}
% RFC-style citations, [RFC3339], linked to the entry. Not \cite: hyperref redefines that.
\newcommand{\xref}[1]{[\hyperref[ref:#1]{#1}]}
\newcommand{\refentry}[2]{\item[{[#1]}]\phantomsection\label{ref:#1}#2}
\newenvironment{refs}{\sloppy\begin{description}[style=nextline,leftmargin=2.5em,font=\normalfont]}{\end{description}}

% A grammar is not split across pages.
\newcounter{rfcfig}
\newcommand{\figcap}[2]{\refstepcounter{rfcfig}\label{#1}{\centering\small Figure~\therfcfig: #2\par}}
\newcommand{\Fig}[1]{Figure~\ref{#1}}

\DefineVerbatimEnvironment{code}{Verbatim}{fontsize=\small,xleftmargin=1.5em,samepage=true}

\begin{document}
\thispagestyle{plain}

% ------------------------------------------------------------------ header block
\begin{tabular*}{\textwidth}{@{}l@{\extracolsep{\fill}}r@{}}
Network Working Group & O. Sigurdsson \\
Internet-Draft & \\
Intended status: Standards Track & \published \\
Expires: \expires & \\
\end{tabular*}

\vspace{2\baselineskip}
\begin{center}
{\Large\bfseries T32: A Compact Universal Timestamp}\\[0.5\baselineskip]
\texttt{\draftname}
\end{center}
\vspace{\baselineskip}

\section*{Abstract}

This document specifies T32, a text and binary encoding of instants. A T32
stamp writes a moment as 8 characters of Crockford base32: a count of
International Atomic Time (TAI) seconds since the start of 2000, offset so
that the past encodes too. 2026-10-31 03:13:00 UTC is \lit{G0S7G6B1}.

A stamp's length is its resolution. Dropping characters names a coarser span
and appending characters names finer time; every stamp names an interval,
never a false instant. Byte-wise sorting is time order, and a prefix is a
range query. The canonical 8-character form covers 15,423 BCE to 19,421 CE at
one-second resolution.

\section*{Note to Readers}

\emph{RFC Editor: please remove this section before publication.}

Source for this draft, its test vectors and two implementations are at
\url{https://github.com/oddurs/t32}. Errata and questions are welcome as
issues there, and proposals in its Discussions. This draft is self-contained:
its requirements are those of draft 0.3 of the T32 specification
\xref{T32}, restated here in full, and a difference between the two is a
defect in one of them. The motivation, terminology, figures and worked example
are informative additions made for this draft. It is prepared for review;
whether to submit it to the IETF is a decision for the project's working group,
not its editor.

\section*{Status of This Memo}

This Internet-Draft is submitted in full conformance with the provisions of
BCP~78 and BCP~79.

Internet-Drafts are working documents of the Internet Engineering Task Force
(IETF). Note that other groups may also distribute working documents as
Internet-Drafts. The list of current Internet-Drafts is at
\url{https://datatracker.ietf.org/drafts/current/}.

Internet-Drafts are draft documents valid for a maximum of six months and may
be updated, replaced, or obsoleted by other documents at any time. It is
inappropriate to use Internet-Drafts as reference material or to cite them
other than as ``work in progress.''

This Internet-Draft will expire on \expires.

\section*{Copyright Notice}

Copyright (c) 2026 IETF Trust and the persons identified as the document
authors. All rights reserved.

This document is subject to BCP~78 and the IETF Trust's Legal Provisions
Relating to IETF Documents (\url{https://trustee.ietf.org/license-info}) in
effect on the date of publication of this document. Please review these
documents carefully, as they describe your rights and restrictions with
respect to this document. Code Components extracted from this document must
include Revised BSD License text as described in Section 4.e of the Trust
Legal Provisions and are provided without warranty as described in the
Revised BSD License.

\clearpage
\tableofcontents
\clearpage

% ================================================================== body

\section{Introduction}\label{sec:intro}

T32 counts TAI seconds, which are uniform, and writes the count in base32 so
that text order is time order. It makes the length of a stamp its resolution:
each character added divides the interval it names by 32. A stamp is then
both a name for an interval and a sort key whose every prefix is a range.
2026-10-31T03:13:00Z is \lit{G0S7G6B1}; \lit{G0S7} names the 12.1 days that
contain it, and \lit{G0S7G6B1G0} the 1/1024 s interval that starts half a
second later.

\subsection{Motivation}\label{sec:motivation}

Each format in use today gives up something T32 aims to keep:

\begin{itemize}
\item \textbf{RFC 3339} \xref{RFC3339}, a profile of ISO~8601 \xref{ISO8601},
is readable but has several spellings of one instant (\lit{Z} or
\lit{+00:00}, any number of fraction digits), so equal instants can differ as
text. It sorts as time only within one spelling and one offset, says how
precise a value is only through how many digits were written, and takes 20
characters for one second.
\item \textbf{POSIX seconds} \xref{POSIX} omit leap seconds: one count repeats
at each insertion, and a difference across one is a second short. A count has
no precision of its own, and its decimal text sorts as time only when
zero-padded.
\item \textbf{TAI64} \xref{TAI64} uses the right time scale and an
offset-binary bias, but its text is hexadecimal at a precision fixed by the
variant, 17 to 33 characters.
\item \textbf{ULID} \xref{ULID} and \textbf{UUIDv7} \xref{RFC9562} sort by
time but are identifiers: a millisecond of POSIX time with random bits
appended, at a resolution that cannot vary.
\end{itemize}

T32 is meant for places where a time is stored, compared or carried as text or
as a fixed-width integer: log records and events; database keys and
partitions, where a prefix is a range; file names and URLs, which want short,
case-insensitive text without punctuation; and data whose times are known only
to the day or the hour, where an instant would overstate them.

\subsection{Goals and Non-Goals}

The goals of T32 are:

\begin{enumerate}
\item \textbf{Uniform.} From 1961 onward, a T32 second is an SI second of TAI,
and subtracting two stamps gives elapsed time. (Before 1961, T32 follows
Universal Time; see \Sec{sec:before1961}.)
\item \textbf{Compact.} 8 characters at one second, against 20 for
\xref{RFC3339} (\lit{2026-10-31T03:13:00Z}).
\item \textbf{Sortable.} Byte-wise order of canonical text is time order.
\item \textbf{Honest resolution.} A stamp's length states the resolution of
the record.
\item \textbf{Hierarchical.} Every prefix is a valid, coarser stamp that
contains the longer one.
\item \textbf{Human-safe.} Case-insensitive, without the most confusable
letters, and with fewer accidental words.
\item \textbf{Stable.} A stamp always denotes the same interval of the T32
time scale. It never needs to change.
\end{enumerate}

The following are not goals:

\begin{itemize}
\item Replacing the civil calendar. People read local dates; T32 sits
underneath them.
\item Time zones. A stamp names an instant everywhere at once.
\item Unique identifiers. T32 carries no randomness. Any random part belongs
in a separate field, or after a character outside the alphabet: appended
alphabet characters would read as resolution.
\end{itemize}

\subsection{Requirements Language}

% The BCP 14 wording is fixed; only its line breaks may give.
{\sloppy The key words ``MUST'', ``MUST NOT'', ``REQUIRED'', ``SHALL'', ``SHALL NOT'',
``SHOULD'', ``SHOULD NOT'', ``RECOMMENDED'', ``NOT RECOMMENDED'', ``MAY'', and
``OPTIONAL'' in this document are to be interpreted as described in
BCP 14 \xref{RFC2119} \xref{RFC8174} when, and only when, they appear in all
capitals, as shown here.\par}

The grammar in \Sec{sec:syntax} uses ABNF \xref{RFC5234}.

\subsection{Terminology}\label{sec:terms}

\begin{description}[style=nextline,leftmargin=1.5em]
\item[Stamp] A sequence of 1 to 16 digits from the T32 alphabet, naming an
interval of the T32 time scale (\Sec{sec:decoding}).
\item[Length ($n$)] The number of digits in a stamp.
\item[Unit ($u(n)$)] The width of the interval a stamp of length $n$ names:
$32^{8-n}$ seconds (\Sec{sec:definitions}).
\item[Value (V)] A stamp read as an unsigned base-32 integer.
\item[Biased time] TAI seconds since the epoch E, plus the bias B
(\Sec{sec:definitions}).
\item[Canonical form] A stamp in uppercase alphabet digits only, with no
hyphens (\Sec{sec:text}).
\item[UTC label] A UTC date and time counted in POSIX seconds
(\Sec{sec:utc}).
\item[\dAT] TAI $-$ UTC as a function of the UTC label (\App{app:dat}).
\item[Inserted and removed time] Instants of TAI that UTC gives no ordinary
label, and labels that name no instant (\Sec{sec:steps}).
\item[Duration] A signed whole number of \pow{2}{-40} s
(\Sec{sec:durations}).
\end{description}

\subsection{Conformance Classes}\label{sec:classes}

This document defines four conformance classes:

\begin{center}
\begin{tabularx}{\textwidth}{@{}l>{\raggedright}p{0.3\textwidth}X@{}}
\toprule
Class & Does & Sections \\
\midrule
Encoder & Instant to stamp & \ref{sec:encoding}, \ref{sec:resolution}, \ref{sec:text}, \ref{sec:binary} \\
Decoder & Stamp to interval & \ref{sec:decoding} \\
Parser & Text to stamp or duration & \ref{sec:syntax}, \ref{sec:text}, \ref{sec:durations} \\
Converter & UTC (or Unix time) to and from TAI & \ref{sec:utc}, \App{app:dat} \\
\bottomrule
\end{tabularx}
\end{center}

An implementation conforms to a class when it meets every requirement of that
class and reproduces every test vector in \App{app:vectors}. The text, those
vectors and \App{app:dat} are normative; where a vector and the text disagree,
the text decides.

\section{Definitions}\label{sec:definitions}

Five constants define T32. Nothing else is configurable.

\begin{description}[style=nextline,leftmargin=1.5em]
\item[Time scale] TAI (International Atomic Time) from 1961; Universal Time
before (see \Sec{sec:utc}). TAI has no leap seconds: every second is equal and
arithmetic is exact.
\item[Epoch E] 2000-01-01T00:00:00 TAI, which is 1999-12-31T23:59:28 UTC and
Unix 946684768. A round, recent instant that keeps present-day values small.
It is not J2000 (2000-01-01T12:00:00 TT).
\item[Bias B] \pow{2}{39} s = 549,755,813,888 s, about 17,421 years. Offset
binary: it centres the range on the epoch, so past instants stay non-negative
and sortable.
\item[Alphabet] \lit{0123456789ABCDEFGHJKMNPQRSTVWXYZ}, Crockford base32
\xref{CROCKFORD}. It omits I, L, O and U, and is in ASCII order, so text order
equals numeric order.
\item[Lengths] 1 to 16 characters; 8 is canonical. Each character carries 5
bits, so 8 characters are 40 bits at 1 s.
\end{description}

A stamp of length $n$ has \textbf{unit} $u(n) = 32^{8-n}$ seconds: 1 s at
length 8, 32 s at length 7, 1/32 s at length 9.

The \textbf{biased time} of an instant $t$ is its TAI seconds since E, plus B.
Every representable instant has a biased time in $[0, 2^{40})$ seconds.

TAI itself is computed after the fact by the BIPM; in real time it is
approximated as UTC(k) + \dAT. T32 does not distinguish these realizations.

\section{Syntax}\label{sec:syntax}

\begin{code}
stamp          = 1*16digit32                  ; canonical form

digit32        = DIGIT / %x41-48 / %x4A-4B / %x4D-4E
                 / %x50-54 / %x56-5A
                 ; 0-9 A-H J K M N P-T V-Z

stamp-input    = in-digit 0*15( [ "-" ] in-digit )

in-digit       = digit32 / lower32 / alias
lower32        = %x61-68 / %x6A-6B / %x6D-6E / %x70-74 / %x76-7A
alias          = %x4F / %x6F / %x49 / %x69 / %x4C / %x6C
                 ; O o read as 0; I i L l read as 1

duration       = sign 1*8digit32 [ "." 1*8digit32 ]
duration-input = sign 1*8in-digit [ "." 1*8in-digit ]
sign           = "+" / "-"

utc-label      = year "-" month "-" day ( "T" / SP )
                 hour ":" minute ":" second [ "." 1*9DIGIT ] "Z"
year           = 4DIGIT / sign 4*6DIGIT
month          = 2DIGIT          ; 01-12
day            = 2DIGIT          ; 01-31, as the month allows
hour           = 2DIGIT          ; 00-23
minute         = 2DIGIT          ; 00-59
second         = 2DIGIT          ; 00-60, 60 only in a leap second
\end{code}

Hyphens in \lit{stamp-input} stand only between two digits, singly. Canonical
durations have further constraints (\Sec{sec:durations}). \lit{utc-label} is
\xref{RFC3339} in UTC, with ISO~8601's expanded years for those outside 0000
to 9999; as in all ABNF, \lit{"T"} and \lit{"Z"} also match lowercase.
Converters read and write UTC labels in this form, with the fewest fraction
digits that are exact.

\section{Encoding}\label{sec:encoding}

An encoder writes instant $t$ at length $n$ as the floor of its biased time
divided by the unit, in $n$ digits:
\[
V = \left\lfloor \frac{t - E + B}{32^{8-n}} \right\rfloor, \qquad 0 \le V < 32^n
\]

\begin{enumerate}
\item Obtain $t - E$ in TAI seconds, exactly (\Sec{sec:utc}).
\item Add B, divide by $u(n)$, and floor. Encoders MUST floor exactly once, at
this step, toward negative infinity; intermediate values MUST NOT be rounded.
Rounding would let a stamp name an interval that does not contain $t$.
\item If V is outside $[0, 32^n)$, the encoder MUST fail. It MUST NOT wrap or
clamp.
\item Encoders MUST write V as exactly $n$ digits from the alphabet, most
significant first, uppercase.
\end{enumerate}

Because there is one floor, encoding at length $n$ and truncating to length
$k$ gives the stamp that encoding at length $k$ gives. Reducing resolution is
a string operation.

\section{Decoding and Interval Semantics}\label{sec:decoding}

A stamp of length $n$ and value V denotes the half-open interval of the T32
time scale (TAI from 1961):
\[
\bigl[\, E - B + V \cdot 32^{8-n},\;\; E - B + (V+1) \cdot 32^{8-n} \,\bigr)
\]

\begin{itemize}
\item \textbf{Start} is the earliest instant the stamp can mean. Sorting and
range queries use it.
\item \textbf{Midpoint} is the best single estimate when one instant is
required.
\item \textbf{Containment.} Stamp A contains stamp B exactly when A is a prefix
of B.
\item \textbf{No partial overlap.} Two stamps are nested or disjoint; there is
no third case.
\item \textbf{Trailing zeros carry meaning.} \lit{G0S7G6B1} and
\lit{G0S7G6B10} are different stamps; the second is one thirty-second of the
first.
\item \textbf{Nanosecond readings.} A decoder that writes an instant as a
nanosecond UTC label MUST write the latest nanosecond label whose length-16
stamp is no later than the instant's length-16 stamp; within a leap second,
labels 23:59:60.f count as labels. Every nanosecond label then survives a
round trip exactly. A label's instant need not be a multiple of \pow{2}{-40}
s, so the instant encoded from a label is its floor, which can fall in time
inserted just before that label (\Sec{sec:steps}); it still reads back as the
label.
\end{itemize}

Each digit selects one of 32 equal parts of the interval named by the digits
before it. \Fig{fig:nest} shows the fifth digit of \lit{G0S7G6B1} selecting
part \lit{G}, the 17th, of \lit{G0S7}.

\begin{code}
G0S7, in 32 parts; the 17th, G, is G0S7G:

 0 1 2 3 4 5 6 7 8 9 A B C D E F G H J K M N P Q R S T V W X Y Z
+-+-+-+-+-+-+-+-+-+-+-+-+-+-+-+-+-+-+-+-+-+-+-+-+-+-+-+-+-+-+-+-+
| | | | | | | | | | | | | | | | |#| | | | | | | | | | | | | | | |
+-+-+-+-+-+-+-+-+-+-+-+-+-+-+-+-+-+-+-+-+-+-+-+-+-+-+-+-+-+-+-+-+
2026-10-24T23:46:35Z                         2026-11-06T03:02:51Z
                                 |
                                 v
G0S7G: 2026-10-31T01:24:43Z to 2026-10-31T10:30:51Z
\end{code}
\figcap{fig:nest}{G0S7G is the 17th of the 32 parts of G0S7}

\section{Ordering and Prefix Properties}\label{sec:ordering}

Byte-wise comparison of canonical stamps orders them by start, and on equal
starts puts the shorter (containing) stamp first.

\begin{itemize}
\item \textbf{Same length:} text order is time order.
\item \textbf{Mixed lengths:} still ordered by start. \lit{G0S7} sorts
immediately before everything it contains.
\item \textbf{Collation:} comparisons MUST be byte-wise (ASCII, or SQL
\lit{COLLATE "C"}). Locale-aware collation can break the order.
\item \textbf{Range query:} everything inside \lit{G0S7}, at any length, is
\lit{t >= \textquotesingle G0S7\textquotesingle{} AND t < \textquotesingle G0S8\textquotesingle}. The upper bound is the prefix with its last
digit incremented, carrying left and dropping every digit that carries to zero
(\lit{G0SZ} becomes \lit{G0T}, \lit{G0ZZ} becomes \lit{G1}); keeping them
would admit the shorter stamp \lit{G0T}, which \lit{G0SZ} does not contain. A
prefix of all \lit{Z} has no upper bound.
\end{itemize}

\section{Length and Resolution}\label{sec:resolution}

Length records the resolution of the record: the width of the interval in
which the instant is stated to lie. It does not describe the accuracy of the
clock that produced it.

\begin{center}
\begin{tabular}{@{}rllll@{}}
\toprule
Length & Unit (exact) & Unit (approx.) & Example & Typical source \\
\midrule
1 & \pow{2}{35} s & 1,089 years & \lit{G} & \\
2 & \pow{2}{30} s & 34.0 years & \lit{G0} & \\
3 & \pow{2}{25} s & 388 days & \lit{G0S} & \\
4 & \pow{2}{20} s & 12.1 days & \lit{G0S7} & \\
5 & \pow{2}{15} s & 9.10 hours & \lit{G0S7G} & \\
6 & \pow{2}{10} s & 17.1 minutes & \lit{G0S7G6} & \\
7 & \pow{2}{5} s & 32 seconds & \lit{G0S7G6B} & \\
\textbf{8} & \textbf{1 s} & \textbf{1 second} & \lit{\textbf{G0S7G6B1}} & \textbf{Seconds (canonical)} \\
10 & \pow{2}{-10} s & 0.977 ms & & Milliseconds \\
12 & \pow{2}{-20} s & 0.954 \textmu s & & Microseconds \\
14 & \pow{2}{-30} s & 0.931 ns & & Nanoseconds \\
16 & \pow{2}{-40} s & 0.909 ps & & Picoseconds \\
\bottomrule
\end{tabular}
\end{center}

Odd lengths are valid. A source's resolution $r$ is the unit in which it
reports time, such as 1 ns for a POSIX \lit{timespec}, not the accuracy of the
clock behind it. An encoder MUST NOT emit a stamp longer than the longest
length $n$ with $u(n) > r/\sqrt{32}$: the length whose unit is nearest $r$ on
a logarithmic scale, the shorter where two are equally near. In integers, $n$
is the longest length with $32 \cdot u(n)^2 > r^2$. A millisecond clock
therefore yields at most 10 characters, a nanosecond clock at most 14.
Encoders SHOULD use one fixed length per field.

\section{Durations}\label{sec:durations}

A duration is a signed whole number of \pow{2}{-40} s, written exactly as a
sign, base32 seconds, and optional base32 fraction digits: \lit{+2MC0} is one
day.

\begin{center}
\begin{tabular}{@{}lrl@{}}
\toprule
Duration & Seconds & T32 duration \\
\midrule
1 minute & 60 & \lit{+1W} \\
1 hour & 3,600 & \lit{+3GG} \\
1 day & 86,400 & \lit{+2MC0} \\
1 week & 604,800 & \lit{+JEM0} \\
Half a second & 0.5 & \lit{+0.G} \\
\bottomrule
\end{tabular}
\end{center}

\begin{itemize}
\item \textbf{Value.} sign $\times$ (I + F / $32^k$) seconds, where I is the
integer part and F the $k$ fraction digits. Every duration is exact in this
form, and its magnitude is less than \pow{2}{40} s.
\item \textbf{Canonical form.} The sign is REQUIRED. No leading zeros in the
integer part (zero is \lit{0}), no trailing zeros in the fraction, no fraction
if it is zero. Zero is \lit{+0}, never \lit{-0}.
\item \textbf{From other units.} Converting a duration from decimal units
(nanoseconds, say) MUST truncate toward zero, so negating a duration negates
its text.
\item \textbf{Stamp minus stamp.} B $-$ A is start(B) $-$ start(A), exactly.
The difference of the instants that B and A stand for lies within
$(-u(A), +u(B))$ of it.
\item \textbf{Range.} Implementations MUST reject a duration whose magnitude
reaches \pow{2}{40} s.
\item \textbf{Stamp plus duration.} Add the duration to the stamp's start,
then floor to the stamp's own length. If the result leaves $[0, 2^{40})$ s,
the operation MUST fail.
\item Durations are SI seconds. \lit{+2MC0} is 86,400 s even across a leap
second; a civil day is a calendar notion.
\end{itemize}

\section{Conversion from UTC}\label{sec:utc}

T32 is TAI-native. UTC converts through \dAT{} = TAI $-$ UTC, given in
\App{app:dat} as a function of the UTC label. \dAT{} has been 37 s since
2017-01-01.

A \textbf{UTC label} L is a time written in the proleptic Gregorian calendar
with astronomical year numbering (year 0 is 1 BCE), counted in seconds since
1970-01-01T00:00:00 by the POSIX formula \xref{POSIX}: days $\times$ 86,400 +
seconds of the day. A Unix timestamp is exactly such a label. Then:
\[
t - E = L + \Delta\mathrm{AT}(L) - 946684800 \text{ s}
\]

Converters MUST evaluate \dAT{} exactly (every value in \App{app:dat} is an
exact decimal) and MUST NOT round before the encoder's single floor.

\subsection{Steps in \texorpdfstring{$\Delta$AT}{Delta AT}}\label{sec:steps}

At each boundary $b$ in \App{app:dat}, \dAT{} changes by
$s = \Delta\mathrm{AT}_\text{after}(b) - \Delta\mathrm{AT}_\text{before}(b)$.
\App{app:dat} lists every non-zero step.

\begin{itemize}
\item \textbf{Removed time ($s < 0$).} A label before $b$ whose TAI instant,
under the old formula, would reach the TAI instant of $b$ under the new one
never occurred. Converters MUST reject it. This happened on 1961-08-01 and
1968-02-01, removing just under 0.05 s and 0.1 s of labels (from
23:59:59.950000001 and 23:59:59.900000003 respectively, at nanosecond
resolution). Defined this way, every instant of TAI still has a label. A
negative leap second would remove 23:59:59.
\item \textbf{Inserted time ($s > 0$).} For $s$ seconds of TAI before the
instant labelled $b$, the formula gives no label.
\begin{itemize}
\item From 1972-07-01 each insertion is one leap second, written 23:59:60 on
the day before $b$. Converters MUST accept second 60 then and only then, and
convert 23:59:60.f to the TAI instant one second after that of 23:59:59.f.
(The insertion of 0.107758 s at 1972-01-01 is fractional and has no second
60.)
\item Decoders MUST write an instant in a leap second as 23:59:60.f, and any
other inserted instant as the label $b$. They MUST report that the instant was
inserted.
\item Before 1972, UTC clocks repeated fractions of a second at these steps;
T32 does not represent the repetition.
\end{itemize}
\end{itemize}

\subsection{Unix Time}

A Unix value is a UTC label and never names 23:59:60. The Unix value at a
leap-second boundary therefore converts to 00:00:00 of the following day. By
the decoder rule, an instant inside a leap second converts back to that same
value.

\subsection{Before 1961}\label{sec:before1961}

Before 1961-01-01, \dAT{} is defined as 0 and the label is read as Universal
Time (mean solar time at Greenwich), extended proleptically. Labels are
Gregorian: a date in the Julian calendar is converted to Gregorian before it is
a label.

\begin{itemize}
\item Pre-1961 stamps are \textbf{stable}: they do not change when models of
Earth's rotation are revised.
\item Pre-1961 T32 time follows Earth's rotation, not atomic seconds. Elapsed
time between such stamps differs from physical time by the change in
$\Delta$T (TT $-$ UT): under about a minute across the last three centuries,
hours in antiquity. Applications that need physical elapsed time in the deep
past apply a $\Delta$T model when they read a stamp, never when they write
one.
\item The definition changes at 1961-01-01, where \dAT{} steps from 0 to
1.4228180 s; the step is inserted time.
\end{itemize}

\subsection{The Table and the Future}\label{sec:future}

In Resolution 5 of October 2026, the 28th CGPM decided that UTC becomes
continuous on 2027-05-20, allowing $|$UT1 $-$ UTC$|$ to reach 3,600 s
\xref{CGPM2026}. IERS
Bulletin C 72 announced no leap second for 2026-12-31 \xref{IERS-C}. No
further step in \dAT{} is therefore expected: \App{app:dat} is expected to be
final, with \dAT{} = 37 s for every label from 2017-01-01.

\begin{itemize}
\item A table carries an expiry date: \App{app:dat} expires at
2027-06-28T00:00:00Z, as does the \lit{leap-seconds.list} it reproduces.
Converting a label at or after the expiry of the table in use, in either
direction, is \textbf{provisional}. A converter MUST return that status with
the result, and MUST NOT fail a conversion for it alone. The status belongs to
the conversion, not to the stamp.
\item Converters SHOULD be able to report which table they used, identified by
its expiry date.
\item If UTC is ever stepped again, \App{app:dat} gains rows in a new version
of this document. Stamps already issued keep their meaning; only conversions
of later labels change.
\end{itemize}

\subsection{TAI Clocks, Smeared Clocks and Stepped Clocks}

\begin{itemize}
\item A converter MAY read TAI directly, for example Linux \lit{CLOCK\_TAI},
where $t - E$ = \lit{CLOCK\_TAI} $-$ 946684800 s. It MUST first confirm that
the kernel's TAI offset equals \dAT{} from its own table (on Linux,
\lit{adjtimex()} reports it as \lit{tai}, which must be 37). An unset offset
is 0, and a stale one is wrong by a whole second or more.
\item Around a leap second, some clocks smear it across many hours and others
repeat a second. Readings from those windows can be up to 1 s from the label
they claim. An encoder SHOULD undo a known smear before encoding. Otherwise,
stamps from the window are not guaranteed to contain the true instant, and
SHOULD be marked as such by the application.
\end{itemize}

\section{Text Form, Display and Parsing}\label{sec:text}

Parse leniently, emit strictly.

\begin{itemize}
\item \textbf{Canonical form.} Encoders MUST emit \lit{stamp}: uppercase
alphabet digits only. Only the canonical form may be stored, compared, hashed
or indexed.
\item \textbf{Display form.} For people, a stamp MAY be shown in groups of
four joined by hyphens: \lit{G0S7-G6B1}, \lit{G0S7-G6B1-3XYE}. The display
form MUST NOT be stored.
\item \textbf{Parsing.} Parsers MUST accept exactly \lit{stamp-input} and MUST
produce the canonical form.
\begin{itemize}
\item Input is ASCII. Any other character, including whitespace and non-ASCII
letters, MUST be rejected. Callers that accept surrounding whitespace strip
it before parsing.
\item Letters are case-insensitive. \lit{O} is read as \lit{0}, and \lit{I} or
\lit{L} as \lit{1}. \lit{U} MUST be rejected.
\item A hyphen MAY appear only between two digits. Leading, trailing and
doubled hyphens MUST be rejected, so a negative duration such as \lit{-1W} is
never read as a stamp.
\item More than 16 digits MUST be rejected.
\item A parser reads its input from left to right and reports the first error
it meets: at a character, the error that character causes; at the end of the
input, a missing sign, a missing part (\lit{empty}) or a trailing hyphen.
Converters check a UTC label in this order: \lit{invalid\_date},
\lit{not\_a\_leap\_second}, \lit{removed\_time}, \lit{out\_of\_range}.
The error codes are listed in \App{app:vectors}.
\end{itemize}
\item \textbf{No check character.} Length carries meaning, so a trailing check
digit would be indistinguishable from one more digit of resolution. Integrity
belongs to the transport.
\end{itemize}

\section{Binary Form}\label{sec:binary}

In binary, a stamp is its value V as an unsigned integer. The length is not
stored: a binary field declares it, as type \textbf{T32-$n$}, and every value
in the field has that length.

\begin{center}
\begin{tabular}{@{}lrrlll@{}}
\toprule
Field type & Bits & Packed bytes & Fits in & Resolution & Typical source \\
\midrule
T32-8 & 40 & 5 & 64-bit integer & 1 s & Seconds \\
T32-10 & 50 & 7 & 64-bit integer & 0.977 ms & Milliseconds \\
T32-12 & 60 & 8 & 64-bit integer & 0.954 \textmu s & Microseconds \\
T32-14 & 70 & 9 & 128-bit integer & 0.931 ns & Nanoseconds \\
T32-16 & 80 & 10 & 128-bit integer & 0.909 ps & Picoseconds \\
\bottomrule
\end{tabular}
\end{center}

\begin{itemize}
\item Packed, V is written big-endian in the fewest whole bytes, and byte
order is time order, like the text. Lengths 8 and 16 pack with no spare bits.
\item As an integer, V sorts in time order within the field.
\item A field SHOULD declare the length that matches its source, by the rule
in \Sec{sec:resolution}. Lengths up to 12 fit a signed 64-bit integer, so
T32-10 suits millisecond sources and T32-12 microsecond sources.
\end{itemize}

\Fig{fig:bits} shows \lit{G0S7G6B1} as digits, bits and a packed T32-8
field. At length 10, \lit{G0S7G6B1G0} is 50 bits; packed as T32-10 in 7 bytes,
its first 6 bits are zero: \lit{0x02 0x00 0xC9 0xE0 0x65 0x86 0x00}.

\begin{code}
Text          G     0     S     7     G     6     B     1
Digit        16     0    25     7    16     6    11     1
Bits      10000 00000 11001 00111 10000 00110 01011 00001
T32-8     10000000 00110010 01111000 00011001 01100001
          0x80     0x32     0x78     0x19     0x61
\end{code}
\figcap{fig:bits}{G0S7G6B1 as text, digit values, bits and packed bytes}

\section{Stability}\label{sec:stability}

This document restates draft 0.3 of the T32 specification \xref{T32}. Until
version 1.0, any part may change. From 1.0, Definitions
(\Sec{sec:definitions}), Encoding (\Sec{sec:encoding}), Decoding
(\Sec{sec:decoding}), the grammar (\Sec{sec:syntax}) and the reading of
labels before 1961 (\Sec{sec:before1961}) are frozen; later versions may only
add, for example rows in \App{app:dat}. A stamp carries no version, and its
meaning never changes.

\section{Implementation Status}\label{sec:impl}

\emph{RFC Editor: please remove this section before publication.}

This section records the status of known implementations at the time of
posting of this Internet-Draft, following \xref{RFC7942}. The description of
implementations here is intended to assist the IETF in its decision processes
in progressing drafts to RFCs. Listing an implementation here does not imply
endorsement by the IETF. The current list is maintained at
\url{https://github.com/oddurs/t32/blob/main/IMPLEMENTATIONS.md}.

\begin{description}[style=nextline,leftmargin=1.5em]
\item[\lit{t32} (Rust)] The reference implementation, published as the
\lit{t32} crate. All four classes; passes every vector from version 0.3.0. Integer arithmetic
only, no required dependencies. Maintained by the author.
\item[\lit{t32.ts} (TypeScript)] The same algorithm in one file, using BigInt.
All four classes; passes every vector. Maintained by the author.
\item[\lit{oracle.py} (Python)] An independent check in exact rational
arithmetic, written from the text and sharing no code with the
implementations. It generates the test vectors. Encoder, decoder, converter
and parser.
\end{description}

Both libraries hold biased time as fixed point in units of \pow{2}{-40} s, and
hold TAI $-$ UTC as exact multiples of $1/(10^7 \times 86{,}400 \times 10^9)$
s, the finest step any coefficient in \App{app:dat} needs. Converting a label
is then exact, and inverting it is a single exact division per period.

No implementation independent of the author exists yet.

\section{Security Considerations}\label{sec:security}

\begin{itemize}
\item \textbf{Canonicalize first.} Different input texts can denote the same
stamp (\lit{g0s7-g6bi}, \lit{G0S7G6B1}). Compare, hash and use as keys only
canonical forms.
\item \textbf{Bounded input.} A valid input is at most 31 characters (16
digits and 15 hyphens), and a duration at most 18. A parser that reports the
first error it meets has met one by the 32nd character, so it never reads
further.
\item \textbf{Confusables.} Parsing is ASCII-only, so look-alike Unicode
letters are rejected rather than mapped.
\item \textbf{Table integrity.} A tampered leap-second table shifts every
conversion after the change. Obtain it from tzdata or the IERS, and check its
expiry date and hash (\lit{leap-seconds.list} carries one).
\item \textbf{Range.} Out-of-range instants and overflowing arithmetic MUST
fail. They MUST NOT wrap.
\item \textbf{No authenticity.} A stamp is a claim about time, as trustworthy
as the clock that produced it and the channel that carried it. T32 provides no
integrity or authenticity. Protocols that use time in security decisions, such
as expiry or the order of authorization events, need to protect a stamp as
they would any other field.
\end{itemize}

\section{Privacy Considerations}\label{sec:privacy}

Resolution is information. A long stamp can reveal timing, such as when an
action happened to the microsecond. Truncate a stamp to the resolution you
mean to publish; for T32 that is a string slice.

\section{IANA Considerations}\label{sec:iana}

This document has no IANA actions. Registrations that could follow, such as a
suffix tag under \xref{RFC9557} or a URN namespace, are listed as an open issue
in \App{app:open}.

% ================================================================== references

\section{References}

\subsection{Normative References}

\begin{refs}
\refentry{IERS-C}{IERS, ``Bulletin C'' (announcements of leap seconds), and
\lit{leap-seconds.list} as distributed with the IANA Time Zone Database,
\url{https://www.iers.org/IERS/EN/Publications/Bulletins/bulletins.html}.}
\refentry{POSIX}{IEEE, ``IEEE Standard for Information Technology---Portable
Operating System Interface (POSIX) Base Specifications, Issue 8'', IEEE Std
1003.1-2024, Base Definitions, ``Seconds Since the Epoch''.}
\refentry{RFC2119}{Bradner, S., ``Key words for use in RFCs to Indicate
Requirement Levels'', BCP~14, RFC~2119, DOI~10.17487/RFC2119, March 1997.}
\refentry{RFC5234}{Crocker, D., Ed. and P. Overell, ``Augmented BNF for Syntax
Specifications: ABNF'', STD~68, RFC~5234, DOI~10.17487/RFC5234, January 2008.}
\refentry{RFC8174}{Leiba, B., ``Ambiguity of Uppercase vs Lowercase in RFC
2119 Key Words'', BCP~14, RFC~8174, DOI~10.17487/RFC8174, May 2017.}
\refentry{TAI-UTC}{U.S. Naval Observatory, \lit{tai-utc.dat}: TAI $-$ UTC
from 1961, \url{https://maia.usno.navy.mil/ser7/tai-utc.dat}, as also embodied
in the IAU SOFA routine \lit{iauDat} (\url{https://www.iausofa.org/}).}
\end{refs}

\subsection{Informative References}

\begin{refs}
\refentry{CGPM2026}{28th General Conference on Weights and Measures (CGPM),
Resolution 5, ``On the technical actions needed to ensure the continuity of
UTC'', October 2026, \url{https://www.bipm.org/en/cgpm-2026/resolution-5}.}
\refentry{CROCKFORD}{Crockford, D., ``Base 32'',
\url{https://www.crockford.com/base32.html}.}
\refentry{ISO8601}{International Organization for Standardization, ``Date and
time --- Representations for information interchange --- Part 1: Basic
rules'', ISO 8601-1:2019.}
\refentry{RFC3339}{Klyne, G. and C. Newman, ``Date and Time on the Internet:
Timestamps'', RFC~3339, DOI~10.17487/RFC3339, July 2002.}
\refentry{RFC7942}{Sheffer, Y. and A. Farrel, ``Improving Awareness of Running
Code: The Implementation Status Section'', BCP~205, RFC~7942,
DOI~10.17487/RFC7942, July 2016.}
\refentry{RFC9557}{Sharma, U. and C. Bormann, ``Date and Time on the Internet:
Timestamps with Additional Information'', RFC~9557, DOI~10.17487/RFC9557,
April 2024.}
\refentry{RFC9562}{Davis, K., Peabody, B., and P. Leach, ``Universally Unique
IDentifiers (UUIDs)'', RFC~9562, DOI~10.17487/RFC9562, May 2024.}
\refentry{T32}{Sigurdsson, O., ``T32: A Compact Universal Timestamp'', draft
0.3, 1 November 2026, tag \lit{spec-draft-0.3},
\url{https://github.com/oddurs/t32/blob/spec-draft-0.3/SPEC.md}.}
\refentry{TAI64}{Bernstein, D. J., ``TAI64, TAI64N, and TAI64NA'',
\url{https://cr.yp.to/libtai/tai64.html}.}
\refentry{ULID}{``Universally Unique Lexicographically Sortable Identifier'',
\url{https://github.com/ulid/spec}.}
\end{refs}

% ================================================================== appendices

\appendix
\makeatletter
\renewcommand{\@seccntformat}[1]{\ifx\csname the#1\endcsname\thesection Appendix \fi\csname the#1\endcsname.\quad}
\makeatother

\section{TAI \texorpdfstring{$-$}{-} UTC}\label{app:dat}

This appendix is normative. \dAT{} = TAI $-$ UTC as a function of the UTC
label, from \xref{IERS-C} and \xref{TAI-UTC}; it expires at
2027-06-28T00:00:00Z. It has been 37 s since
2017-01-01, and with UTC continuous from 2027-05-20 it is expected to stay
37 s (\Sec{sec:future}). Before 1961-01-01, \dAT{} is defined as 0 for T32.

\textbf{1972 onward} (whole seconds; each value applies from 00:00:00 UTC on
its date):

\begin{center}
\begin{tabular}{@{}lr@{\qquad}lr@{}}
\toprule
From & \dAT{} (s) & From & \dAT{} (s) \\
\midrule
1972-01-01 & 10 & 1985-07-01 & 23 \\
1972-07-01 & 11 & 1988-01-01 & 24 \\
1973-01-01 & 12 & 1990-01-01 & 25 \\
1974-01-01 & 13 & 1991-01-01 & 26 \\
1975-01-01 & 14 & 1992-07-01 & 27 \\
1976-01-01 & 15 & 1993-07-01 & 28 \\
1977-01-01 & 16 & 1994-07-01 & 29 \\
1978-01-01 & 17 & 1996-01-01 & 30 \\
1979-01-01 & 18 & 1997-07-01 & 31 \\
1980-01-01 & 19 & 1999-01-01 & 32 \\
1981-07-01 & 20 & 2006-01-01 & 33 \\
1982-07-01 & 21 & 2009-01-01 & 34 \\
1983-07-01 & 22 & 2012-07-01 & 35 \\
 & & 2015-07-01 & 36 \\
 & & 2017-01-01 & 37 \\
\bottomrule
\end{tabular}
\end{center}

\textbf{1961 to 1971} (\dAT{} = $a + (\mathrm{MJD} - m) \times r$ seconds,
where MJD is the Modified Julian Date of the UTC label, fraction of the day
included):

\begin{center}
\begin{tabular}{@{}lrrr@{}}
\toprule
From & $a$ (s) & $m$ (MJD) & $r$ (s/day) \\
\midrule
1961-01-01 & 1.4228180 & 37300 & 0.0012960 \\
1961-08-01 & 1.3728180 & 37300 & 0.0012960 \\
1962-01-01 & 1.8458580 & 37665 & 0.0011232 \\
1963-11-01 & 1.9458580 & 37665 & 0.0011232 \\
1964-01-01 & 3.2401300 & 38761 & 0.0012960 \\
1964-04-01 & 3.3401300 & 38761 & 0.0012960 \\
1964-09-01 & 3.4401300 & 38761 & 0.0012960 \\
1965-01-01 & 3.5401300 & 38761 & 0.0012960 \\
1965-03-01 & 3.6401300 & 38761 & 0.0012960 \\
1965-07-01 & 3.7401300 & 38761 & 0.0012960 \\
1965-09-01 & 3.8401300 & 38761 & 0.0012960 \\
1966-01-01 & 4.3131700 & 39126 & 0.0025920 \\
1968-02-01 & 4.2131700 & 39126 & 0.0025920 \\
\bottomrule
\end{tabular}
\end{center}

\textbf{Steps} (from these tables):

\begin{center}
\begin{tabularx}{\textwidth}{@{}Xrl@{}}
\toprule
Boundary & Step (s) & Kind \\
\midrule
1961-01-01 & +1.4228180 & Inserted (start of atomic definition) \\
1961-08-01 & $-$0.05 & Removed \\
1963-11-01, 1964-04-01, 1964-09-01, 1965-01-01, 1965-03-01, 1965-07-01, 1965-09-01 & +0.1 each & Inserted \\
1968-02-01 & $-$0.1 & Removed \\
1972-01-01 & +0.107758 & Inserted \\
Each date from 1972-07-01 in the first table & +1 & Inserted (leap second, 23:59:60) \\
\bottomrule
\end{tabularx}
\end{center}

The boundaries 1962-01-01, 1964-01-01 and 1966-01-01 change the formula but not
its value.

\section{Test Vectors}\label{app:vectors}

This appendix is normative: a conforming implementation reproduces every
vector in it for its class. A larger set of 782 vectors, including both sides
of every boundary in \App{app:dat}, is published with \xref{T32} as
\lit{vectors/t32-vectors.json} at the tag \lit{spec-draft-0.3}. It is
informative here, because an Internet-Draft carries its normative content
itself; an implementation that passes it passes the vectors below.

Every vector was computed by an exact rational-arithmetic oracle written
directly from the text (\Sec{sec:impl}), errors included, and is reproduced
by both libraries. A disagreement between this text, the vectors and an
implementation is a defect.

Errors in the vectors file are named by code:

\begin{center}
\begin{tabularx}{\textwidth}{@{}lX@{}}
\toprule
Code & Condition \\
\midrule
\lit{empty} & No digit where one is required: the whole stamp, or either part of a duration \\
\lit{too\_long} & More than 16 digits in a stamp, or 8 in either part of a duration \\
\lit{invalid\_char} & A character outside the input grammar \\
\lit{misplaced\_hyphen} & A hyphen not between two digits \\
\lit{missing\_sign} & A duration that does not begin with \lit{+} or \lit{-} \\
\lit{invalid\_date} & Text that is not a \lit{utc-label}, or a day or time the calendar lacks \\
\lit{not\_a\_leap\_second} & Second 60 where no leap second was inserted \\
\lit{removed\_time} & A UTC label in removed time \\
\lit{out\_of\_range} & An instant, stamp or duration outside its range \\
\bottomrule
\end{tabularx}
\end{center}

\subsection{Encoder and Converter (UTC Labels)}

{\small
\begin{longtable}{@{}llrl@{}}
\toprule
Case & UTC & Length & Stamp \\
\midrule
\endhead
T32 epoch & 1999-12-31T23:59:28 & 8 & \lit{G0000000} \\
Unix epoch (\dAT{} = 8.000082 s) & 1970-01-01T00:00:00 & 8 & \lit{FZ3S5F48} \\
Apollo 11 landing & 1969-07-20T20:17:40 & 8 & \lit{FZ3BMMKB} \\
Rubber era, exact (ns clock) & 1966-06-15T12:00:00.123456789 & 14 & \lit{FZ0EDKA4VPC2A1} \\
UTC becomes integer & 1972-01-01T00:00:00 & 8 & \lit{FZ5NA8WA} \\
Before last leap second & 2016-12-31T23:59:59 & 8 & \lit{G0FZP0S3} \\
Last leap second & 2016-12-31T23:59:60 & 8 & \lit{G0FZP0S4} \\
After last leap second & 2017-01-01T00:00:00 & 8 & \lit{G0FZP0S5} \\
Unix 1483228800 & (2017-01-01T00:00:00) & 8 & \lit{G0FZP0S5} \\
Spec date & 2026-10-31T03:13:00 & 8 & \lit{G0S7G6B1} \\
Spec date + 0.5 s (ms clock) & 2026-10-31T03:13:00.500 & 10 & \lit{G0S7G6B1G0} \\
Spec date + 0.123 s (\textmu s clock) & 2026-10-31T03:13:00.123000 & 12 & \lit{G0S7G6B13XYE} \\
Before 1961, taken as UT & 1675-08-20T00:00:00 & 8 & \lit{FPEY6N00} \\
\bottomrule
\end{longtable}
}

\subsection{Converter Errors}

\begin{center}
\begin{tabular}{@{}ll@{}}
\toprule
UTC & Result \\
\midrule
1961-07-31T23:59:59.95 & Accepted \\
1961-07-31T23:59:59.950000001 & Error: removed time \\
1968-01-31T23:59:59.900000002 & Accepted \\
1968-01-31T23:59:59.900000003 & Error: removed time \\
2016-12-30T23:59:60 & Error: not a leap second \\
\bottomrule
\end{tabular}
\end{center}

\subsection{Decoder (Interval Start in TAI Seconds since E)}

\begin{center}
\begin{tabular}{@{}lrr@{}}
\toprule
Stamp & Start (s) & Width (s) \\
\midrule
\lit{00000000} & $-$549,755,813,888 & 1 \\
\lit{FZ3BMMKB} & $-$960,867,733 & 1 \\
\lit{G0000000} & 0 & 1 \\
\lit{G0FZP0S4} & 536,544,036 & 1 \\
\lit{G0S7} & 846,200,832 & 1,048,576 \\
\bottomrule
\end{tabular}
\end{center}

\subsection{Decoder with Converter (Interval as UTC Labels, Floored to 1 ns)}

\begin{center}
\begin{tabularx}{\textwidth}{@{}l>{\raggedright\arraybackslash}Xl@{}}
\toprule
Stamp & Start & End (exclusive) \\
\midrule
\lit{00000000} & $-$15422-12-06T11:41:52 & $-$15422-12-06T11:41:53 \\
\lit{FZ3BMMKB} & 1969-07-20T20:17:39.425406217 & 1969-07-20T20:17:40.425406187 \\
\lit{G0FZP0S4} & 2016-12-31T23:59:60 (inserted; label 2017-01-01T00:00:00) & 2017-01-01T00:00:00 \\
\lit{G0S7} & 2026-10-24T23:46:35 & 2026-11-06T03:02:51 \\
\lit{G0S7G6} & 2026-10-31T03:07:07 & 2026-10-31T03:24:11 \\
\bottomrule
\end{tabularx}
\end{center}

The \lit{FZ3BMMKB} interval is 0.99999997 UTC seconds wide, because UTC seconds
were stretched in 1969. Years are astronomical ($-$15422 is 15423 BCE).

\subsection{Parser}

\begin{center}
\begin{tabular}{@{}ll@{}}
\toprule
Input & Result \\
\midrule
\lit{g0s7-g6bi} & \lit{G0S7G6B1} \\
\lit{oooo} & \lit{0000} \\
\lit{G0S7G6B1-3XYE} & \lit{G0S7G6B13XYE} \\
\lit{G0SU} & Error: invalid character \\
\lit{\textvisiblespace G0S7} (leading space) & Error: invalid character \\
\lit{\i} (U+0131) & Error: invalid character \\
\lit{-1W}, \lit{1W-}, \lit{G0S7-{}-G6B1} & Error: misplaced hyphen \\
17 digits & Error: too long \\
empty & Error: empty \\
\bottomrule
\end{tabular}
\end{center}

\subsection{Durations}

\begin{center}
\begin{tabular}{@{}lll@{}}
\toprule
Input & Value & Canonical \\
\midrule
86,400 s & one day & \lit{+2MC0} \\
$-$1.5 s & & \lit{-1.G} \\
$-$1 ns (truncated toward zero) & $-$1,099 $\times$ \pow{2}{-40} s & \lit{-0.0000012B} \\
0 & & \lit{+0} \\
\lit{G0S8} start $-$ \lit{G0S7} start & & \lit{+10000} \\
\lit{G0S7G6B1} start $-$ \lit{G0S7} start & & \lit{+G6B1} \\
\lit{1W} (no sign) & & Error \\
\bottomrule
\end{tabular}
\end{center}

\section{Worked Example}\label{app:example}

This appendix is informative. It follows one instant through encoding,
truncation, decoding, a range query and a duration.

\textbf{Encoding 2026-10-31T03:13:00Z at length 8} (\Sec{sec:encoding}):

\begin{enumerate}
\item The UTC label L is 1793416380 POSIX seconds.
\item \dAT(L) is 37 s (\App{app:dat}, from 2017-01-01).
\item $t - E = L + \Delta\mathrm{AT}(L) - 946684800$ = 846,731,617 s.
\item Adding B = 549,755,813,888 s gives a biased time of
550,602,545,505 s.
\item The unit at length 8 is 1 s, so V = 550,602,545,505.
\item In base 32, V has the digits 16, 0, 25, 7, 16, 6, 11, 1, written
\lit{G0S7G6B1} (\Fig{fig:bits}).
\end{enumerate}

% Long numbers cannot break; let this paragraph's other lines give instead.
{\sloppy\textbf{A finer stamp.} From a millisecond clock reading 03:13:00.500, the
biased time is 550,602,545,505.5 s. At length 10 the unit is \pow{2}{-10} s,
so
\[
V = \lfloor 550{,}602{,}545{,}505.5 \times 1024 \rfloor = 563{,}817{,}006{,}597{,}632,
\]
written \lit{G0S7G6B1G0}. Its first 8 digits are the length-8 stamp, as the
single floor guarantees.\par}

\textbf{Truncation and decoding.} Each prefix of \lit{G0S7G6B1} is a coarser
stamp that contains it:

\begin{center}
\begin{tabular}{@{}lllr@{}}
\toprule
Stamp & Starts (UTC) & Ends, exclusive (UTC) & Width \\
\midrule
\lit{G0S7} & 2026-10-24T23:46:35 & 2026-11-06T03:02:51 & 12.1 days \\
\lit{G0S7G} & 2026-10-31T01:24:43 & 2026-10-31T10:30:51 & 9.10 hours \\
\lit{G0S7G6} & 2026-10-31T03:07:07 & 2026-10-31T03:24:11 & 17.1 minutes \\
\lit{G0S7G6B} & 2026-10-31T03:12:59 & 2026-10-31T03:13:31 & 32 seconds \\
\lit{G0S7G6B1} & 2026-10-31T03:13:00 & 2026-10-31T03:13:01 & 1 second \\
\bottomrule
\end{tabular}
\end{center}

For \lit{G0S7}, V = 16 $\times$ \pow{32}{3} + 0 $\times$ \pow{32}{2} + 25
$\times$ 32 + 7 = 525,095 and the unit is \pow{2}{20} = 1,048,576 s, so the
interval starts at 525,095 $\times$ 1,048,576 $-$ B = 846,200,832 s after E
(\Sec{sec:decoding}). Subtracting \dAT{} and adding 946684800 gives the UTC
label 1792885595, which is 2026-10-24T23:46:35Z.

\textbf{Range query.} Every stamp inside \lit{G0S7}, at any length, satisfies
\lit{t >= \textquotesingle G0S7\textquotesingle{} AND t < \textquotesingle G0S8\textquotesingle} under byte-wise
comparison (\Sec{sec:ordering}).

\textbf{Duration.} The start of \lit{G0S7G6B1} minus the start of
\lit{G0S7} is 530,785 s exactly, written \lit{+G6B1} (\Sec{sec:durations}).

\section{Prior Art}\label{app:prior}

This appendix is informative.

The example column writes 2026-10-31T03:13:00Z in each format.

\begin{center}
\footnotesize
\begin{tabularx}{\textwidth}{@{}>{\raggedright}p{5.6em}lr>{\raggedright}p{4.6em}>{\raggedright}X>{\raggedright\arraybackslash}X@{}}
\toprule
Format & Example & Chars & Time scale & Text sorts as time & Resolution in the value \\
\midrule
RFC 3339, ISO 8601 & \lit{2026-10-31T03:13:00Z} & 20 & UTC & Within one profile and offset & Partly \\
Unix time & \lit{1793416380} & 10 & UTC, POSIX & When zero-padded & No \\
TAI64 label & \lit{@400000006ae55ce1} & 17 & TAI & Yes & Fixed per variant \\
ULID, time part & \lit{01M67Z4FK0} & 10 & Unix, ms & Yes & No \\
UUIDv7 & 48-bit Unix ms & 36 & Unix, ms & Yes & No \\
\textbf{T32} & \lit{\textbf{G0S7G6B1}} & \textbf{8} & \textbf{TAI} & \textbf{Yes} & \textbf{Yes, by length} \\
\bottomrule
\end{tabularx}
\end{center}

TAI64 \xref{TAI64} is the closest relative: the same timescale and the same
offset-binary idea (its bias is \pow{2}{62}), written in hexadecimal at a
precision fixed by the variant (TAI64, TAI64N, TAI64NA). ULID \xref{ULID}
showed that sortable Crockford base32 works for time-ordered keys, and UUIDv7
\xref{RFC9562} brought a time-ordered layout into a standard identifier. T32
adds three things: length as resolution, with interval semantics and prefix
containment; the shortest case-insensitive text that sorts as time; and an
exact UTC mapping at every step of the table.

\section{Known Limits}\label{app:limits}

This appendix is informative. These limits are deliberate; removing any of
them would cost a goal.

\begin{itemize}
\item \textbf{Calendar alignment.} Buckets do not fall on days, months or
years. Eras are aligned to the epoch: \lit{F} covers about 911 to 2000 and
\lit{G} about 2000 to 3089, so dates either side of 2000 differ in the first
character.
\item \textbf{Relativity.} TAI, like Terrestrial Time (TT = TAI + 32.184 s),
is a time scale for clocks on Earth. Clocks in orbit, on the Moon or on Mars
convert through a defined relativistic transformation, as GPS does, before
encoding.
\item \textbf{Before 1961.} Stamps follow Universal Time, not atomic seconds
(\Sec{sec:before1961}).
\item \textbf{Range.} The canonical form ends in 19,421 CE. Extending it would
add a character to every stamp in use.
\item \textbf{Uniqueness.} Two events in the same unit share a stamp. T32 is a
time, not an identifier; see \Sec{sec:intro} for where randomness may go.
\end{itemize}

\section{Open Issues}\label{app:open}

\emph{RFC Editor: please remove this appendix before publication.}

These are the questions the author considers open. Each states the current
position and the kind of evidence that would change it.

\begin{description}[style=nextline,leftmargin=1.5em]
\item[The name] ``T32'' is also the name of Arm's 32-bit Thumb instruction set
and the common name of Lauterbach's TRACE32 debugger. The current position is
to keep the name, since the contexts rarely meet and the name says what the
format is: T for time, 32 for its radix. Evidence of real confusion, or a
better name, would reopen it.
\item[The epoch] T32 counts from 2000-01-01T00:00:00 TAI. Alternatives
considered were J2000 (2000-01-01T12:00:00 TT), which suits astronomy, and
1972-01-01, when UTC's offset from TAI became whole seconds. The current
position favours a round instant in the scale T32 counts. Since the epoch
fixes every stamp, this question closes before publication.
\item[Before 1961] Labels before 1961 are read as Universal Time, which keeps
stamps stable but makes elapsed time non-physical across centuries. The
alternative is to apply a $\Delta$T model, which is physical but changes
stamps whenever the model is revised. The current position favours stability.
The reading of these labels is frozen at version 1.0, so this question closes
then.
\item[The resolution rule] An encoder may use the length whose unit is nearest
its source's resolution on a logarithmic scale, so a stamp can claim up to
$\sqrt{32}$ (about 5.7) times more resolution than its source has. The
stricter alternative is the longest length whose unit is no finer than the
source, which gives 9 for a millisecond clock instead of 10. The current
position prefers the familiar lengths 8, 10, 12 and 14.
\item[Letter case] Canonical stamps are uppercase. Lowercase reads better in
URLs and file names, and some systems fold case. Since parsing is
case-insensitive, the question is only which form is canonical. The current
position is uppercase, as in Crockford base32.
\item[A human rendering] This document defines a display form for stamps, but
not how to show a stamp's interval to people in their own calendar and time
zone. The question is whether that belongs here or in a companion document.
\item[Identifiers and registration] T32 defines no URI scheme, media type or
registered tag. Candidates include a suffix tag under \xref{RFC9557} and a URN
namespace. Registration would follow only if implementers ask for it.
\item[Other time scales] TAI suits clocks on Earth. Lunar and Martian time
scales are being defined, and a T32 profile for one of them would need its own
epoch and conversion. The current position is to leave the format TAI-only and
to define profiles, if any, in separate documents.
\end{description}

\section*{Author's Address}
\addcontentsline{toc}{section}{Author's Address}

Oddur Sigurdsson\\
Email: \href{mailto:oddurs@gmail.com}{oddurs@gmail.com}

\end{document}
